\documentclass[11pt,a4paper]{article}
\usepackage[T1]{fontenc}
\usepackage[utf8]{inputenc}
\usepackage{lmodern}
\usepackage[a4paper,top=1.9cm,bottom=1.8cm,left=2.3cm,right=2.3cm]{geometry}
\usepackage{amsmath,amssymb}
\usepackage{xcolor}
\usepackage{longtable,array,booktabs}
\usepackage{fancyhdr}
\newcommand{\barre}[1]{\leavevmode\setbox0\hbox{#1}\rlap{\rule[0.55ex]{\wd0}{0.5pt}}\box0}
\definecolor{bleu}{HTML}{1F4E79}
\pagestyle{fancy}\fancyhf{}
\lhead{\small 5\textsuperscript{e} Mathématiques --- Trigonométrie}
\rhead{\small Partie 2 : formules et équations}
\cfoot{\small\thepage}
\rfoot{\footnotesize Aide de l'IA (AIA)}
\renewcommand{\headrulewidth}{0.4pt}
\setlength{\parindent}{0pt}
% ---- lignes du tableau ----
\newcommand{\case}{$\square$}
\newcommand{\sect}[1]{\midrule\multicolumn{2}{@{}l}{\rule{0pt}{4ex}\textbf{\large\color{bleu}#1}}\\[0.6ex]\midrule}
\newcommand{\ligne}[1]{#1 & \case\\}                         % savoir-faire simple
\newcommand{\groupe}[1]{\rule{0pt}{3.6ex}\textbf{#1} & \\} % titre d'un groupe
\newcommand{\sous}[1]{\hangindent=2.1em\hangafter=1\hspace{1.2em}$\circ$ #1 & \case\\}  % ligne d'un groupe
\begin{document}
\begin{center}
{\Large\bfseries\color{bleu} Tableau des savoir-faire}\\[1mm]
Trigonométrie $\bullet$ Partie 2 : formules et équations trigonométriques
\end{center}
\vspace{-2mm}
\renewcommand{\arraystretch}{1.04}
\begin{longtable}{@{}>{\raggedright\arraybackslash}p{12.6cm}>{\centering\arraybackslash}p{2.9cm}@{}}
\toprule
 & \textbf{Éval. (x ou v)}\\
\endhead
\bottomrule
\endlastfoot

\sect{À revoir avant de commencer}
\ligne{Les valeurs particulières de $\sin$, $\cos$ et $\tan$}
\ligne{Les angles associés}
\ligne{La formule fondamentale : $\sin^2a+\cos^2a=1$}

\sect{1. Connaître et démontrer les formules}
\ligne{Connaître les formules d'addition : $\sin(a\pm b)$, $\cos(a\pm b)$, $\tan(a\pm b)$}
\ligne{Connaître les formules de duplication : $\sin(2a)$, $\cos(2a)$, $\tan(2a)$}
\ligne{Connaître les formules de Carnot : $\cos(2a)=1-2\sin^2a=2\cos^2a-1$}
\barre{Connaître les formules de Simpson}\quad\textit{(elles sont données)} & \\
\groupe{Démontrer une formule}
\sous{démontrer la formule de $\sin(a+b)$}
\sous{en déduire les autres formules d'addition ($\sin$, $\cos$ et $\tan$ de $a\pm b$)}
\sous{démontrer les formules de duplication (avec les formules d'addition)}
\sous{démontrer les formules de Carnot (avec $\cos(2a)$ et la formule fondamentale)}
\sous{démontrer les formules de Simpson (avec les formules d'addition)}

\sect{2. Utiliser les formules}
\groupe{Calculer une valeur exacte, sans calculatrice}
\sous{d'une somme ou d'une différence d'angles remarquables (ex. : $\cos 75^\circ$, $\sin\frac{\pi}{12}$)}
\sous{de $\sin(a+b)$, $\cos(a-b)$, $\tan(a+b)$, \dots{} à partir de nombres trigonométriques de $a$ et de $b$}
\sous{de $\sin(2a)$, $\cos(2a)$, $\tan(2a)$ à partir d'un nombre trigonométrique de $a$ et de son quadrant}
\groupe{Transformer une expression}
\sous{écrire $\cos a$ à partir de $\sin a$ (ou l'inverse) avec la formule fondamentale}
\sous{développer ou regrouper avec les formules d'addition}
\sous{avec les formules de duplication}
\sous{avec les formules de Carnot (remplacer $\sin^2a$ ou $\cos^2a$)}
\sous{transformer une somme en produit (Simpson)}
\sous{transformer un produit en somme (Simpson)}
\groupe{Raisonner avec les formules}
\sous{dire si une égalité est vraie ou fausse, et justifier (ex. : $\sin(2a)=2\sin a$)}
\sous{vérifier une identité trigonométrique}

\pagebreak
\sect{3. Résoudre une équation trigonométrique}
\groupe{Résoudre une équation élémentaire}
\sous{$\sin x = c$}
\sous{$\cos x = c$}
\sous{$\tan x = c$ (avec sa condition d'existence)}
\sous{reconnaître une équation sans solution (ex. : $\cos x = 2$)}
\groupe{Écrire les solutions}
\sous{donner toutes les solutions (avec $+\,2k\pi$ ou $+\,k\pi$, $k\in\mathbb{Z}$)}
\sous{donner les solutions qui se trouvent dans un intervalle donné}
\sous{représenter les solutions sur le cercle trigonométrique}
\groupe{Résoudre une équation qui se ramène à une équation élémentaire}
\sous{l'inconnue est dans une expression : $\sin(ax+b)=c$, $\cos(ax+b)=c$, \dots}
\sous{égalité de la même famille : $\sin A=\sin B$, $\cos A=\cos B$, $\tan A=\tan B$}
\sous{égalité de familles différentes : $\sin A=\cos B$ (angles associés)}
\groupe{Résoudre une équation en choisissant la méthode}
\sous{par mise en évidence ou factorisation (produit nul)}
\sous{en se ramenant à une équation du 2\textsuperscript{e} degré (poser $y=\sin x$ ou $y=\cos x$)}
\sous{en utilisant d'abord une formule (duplication, Carnot, Simpson)}
\sous{équation linéaire $a\sin x+b\cos x=c$}
\sous{équation homogène en $\sin x$ et $\cos x$}

\sect{4. Résoudre une inéquation trigonométrique}
\ligne{Résoudre une inéquation élémentaire ($\sin x \geq c$, $\cos x < c$, \dots) à l'aide du cercle trigonométrique}
\end{longtable}
\end{document}
